\section{Metodologia}
\label{sec:metodologia}

\subsection{Delineamento da pesquisa}

A pesquisa é aplicada, de abordagem mista e objetivo exploratório-descritivo.
O eixo de construção seguirá a \textit{Design Science Research Methodology}
(DSRM), que organiza o trabalho em identificação do problema, definição dos
objetivos, projeto e desenvolvimento, demonstração, avaliação e comunicação
\cite{peffersEtAl2007}. A adoção da DSRM é adequada porque a contribuição inclui
um artefato destinado a resolver um problema organizacional, cuja utilidade
precisa ser demonstrada e avaliada \cite{hevnerEtAl2004}.

Em paralelo, será realizado um estudo de caso único, com unidades de análise
incorporadas correspondentes a funções e, quando autorizado, a diferentes
unidades da rede. O estudo de caso contextualizará o processo, os incidentes e
as restrições organizacionais, seguindo recomendações de planejamento,
triangulação e relato de \cite{runesonHost2009}. A empresa será identificada no
texto apenas como ``uma rede de postos de combustíveis do sul do Brasil''.

\begin{table}[H]
  \centering
  \caption{Etapas da DSRM e evidências previstas}
  \label{tab:dsrm}
  \small
  \begin{tabularx}{\textwidth}{@{}L{3.1cm}YL{2.5cm}@{}}
    \toprule
    \textbf{Etapa} & \textbf{Evidência ou produto} & \textbf{Estado} \\
    \midrule
    Identificar problema & Registros anonimizados, entrevistas e descrição do processo & Parcial; dados pendentes \\
    Definir objetivos & Pergunta, hipóteses, objetivos e critérios de aceitação & Elaborado \\
    Projetar e desenvolver & Requisitos, arquitetura, simuladores, software e bancada & Software parcial; bancada pendente \\
    Demonstrar & Roteiro reproduzível com casos compatíveis, incompatíveis e falhas & Componentes simulados; roteiro integral pendente \\
    Avaliar & Testes funcionais, falhas, latência, usabilidade e cenários econômicos & Verificação automatizada parcial \\
    Comunicar & Artigo, repositório, seminário e defesa & Em andamento \\
    \bottomrule
  \end{tabularx}
  \\[2pt]\fontepropria
\end{table}

\subsection{Caracterização do processo e entrevistas}

Serão convidados, conforme autorização e disponibilidade, frentistas de mais de
uma unidade, gestores operacionais, representante da direção e profissional do
setor financeiro. A amostra será intencional e buscará diversidade de papéis,
turnos e unidades, sem pretensão de representatividade estatística. O roteiro
semiestruturado abrangerá sequência de atendimento, pontos de conferência,
causas percebidas, tratamento de exceções, consequências e aceitação de uma
barreira automatizada.

As entrevistas somente ocorrerão após autorização escrita da empresa,
orientação formal sobre o enquadramento ético e, quando aplicável, aprovação do
Comitê de Ética em Pesquisa (CEP). Serão adotados consentimento, participação
voluntária e plano de anonimização, retenção e descarte coerentes com a
Resolução CNS n\textsuperscript{o} 510/2016 \cite{cns5102016}. Gravação será opcional e dependerá
de consentimento específico. A análise utilizará codificação temática; as
categorias, decisões de codificação e trechos anonimizados serão registrados
para auditabilidade, sem expor pessoas ou unidades.

\pendente{registrar a decisão do orientador/CEP, a autorização da empresa, os
critérios finais de recrutamento e o número efetivamente entrevistado.}

\subsection{Dados históricos e custos}

A coleta documental solicitará um conjunto mínimo: identificador artificial da
ocorrência, período, unidade pseudonimizada, categorias de combustível,
momento de detecção, fonte documental e componentes comprovados do custo. Placa,
nome de cliente, nome de funcionário, chassi, Renavam e relato livre
identificável não serão inseridos na base de pesquisa.

Os valores serão classificados em custo documentado, estimativa e impacto
qualitativo. Somente itens suportados por registros serão somados como perdas
realizadas. Quando a série abranger datas diferentes, a correção para a data-base
usará o IPCA oficial \cite{ibgeIpca2026}, conforme a Equação~\ref{eq:ipca}:

\begin{equation}
  V_{corrigido}=V_{evento}\times
  \frac{IPCA_{data\mbox{-}base}}{IPCA_{evento}}.
  \label{eq:ipca}
\end{equation}

A taxa de incidentes por abastecimento somente será calculada se numerador e
denominador tiverem mesma população, período e unidades. Sem essa compatibilidade,
serão apresentadas apenas contagens e custos absolutos. O dicionário e as regras
de qualidade estão em \texttt{research/data/dicionario-dados.md}.

\subsection{Construção do artefato}

O desenvolvimento é incremental. Já foram implementados o domínio, o provedor
veicular simulado, a política de compatibilidade, o simulador de bicos, a API,
a auditoria minimizada em SQLite e uma PWA. Também existe firmware experimental,
cuja lógica portátil foi verificada em computador e cujo alvo \path{esp32dev}
foi compilado com PlatformIO. Autenticação individual,
persistência da versão da política, gateway MQTT no backend e integração entre
os componentes permanecem pendentes. A bancada física somente será montada após
validação do circuito e dos componentes com o orientador. A integração real com
SENATRAN permanecerá desativada até a obtenção das autorizações e credenciais
necessárias.

Cada cenário de demonstração terá entrada, estado inicial, eventos, decisão
esperada, comando esperado e estado final. A política e os contratos de
interface serão mantidos independentes da tecnologia física, permitindo trocar
um atuador simulado por um controlador de bancada sem modificar a regra de
negócio.

\subsection{Plano de avaliação}

A avaliação combinará cinco dimensões:

\begin{enumerate}
  \item \textbf{correção funcional:} decisão para veículos e produtos
  compatíveis, incompatíveis, multicombustível e desconhecidos;
  \item \textbf{segurança diante de falhas:} timeout, resposta inválida, perda
  de comunicação, reinicialização, mensagem duplicada e expiração;
  \item \textbf{desempenho:} latência ponta a ponta com mediana, percentil 95 e
  máximo, separando provedor, backend e dispositivo;
  \item \textbf{usabilidade:} conclusão da tarefa, erros, tempo e percepção dos
  participantes autorizados, com estatística descritiva; e
  \item \textbf{viabilidade econômica:} custo anualizado e perdas evitáveis sob
  cenários de 25\%, 50\%, 75\% e 100\% de prevenção.
\end{enumerate}

Sejam $Q$ o número anual de consultas, $c_q$ o custo unitário, $C_{GCC}$ a
remuneração anual da GCC, $C_{infra}$ infraestrutura, $C_{man}$ manutenção e
$C_{hw,a}$ a parcela anualizada do hardware. O custo anual da solução será:

\begin{equation}
  C_{\mathrm{sol}}=Q c_q+C_{GCC}+C_{infra}+C_{man}+C_{hw,a}.
  \label{eq:custo-solucao}
\end{equation}

Para uma perda histórica anual corrigida $L$ e percentual de prevenção $p$, o
benefício líquido contrafactual será $B(p)=pL-C_{\mathrm{sol}}$. Quando $L>0$, o ponto
de equilíbrio será $p^{*}=C_{\mathrm{sol}}/L$. O modelo não incluirá ganho de imagem,
fidelização ou risco jurídico sem método e evidência próprios.

\subsection{Ameaças à validade}

As principais ameaças previstas são registros incompletos, memória seletiva nas
entrevistas, pequena amostra, influência do pesquisador que também trabalha na
empresa, uso de simulador, diferença entre bancada e pista e mudanças de preço
ou regulação. As estratégias de mitigação incluem triangulação entre fontes,
roteiro comum, registro de decisões, dados sintéticos versionados, repetição dos
ensaios, relato de casos divergentes, revisão pelo orientador e atualização das
fontes oficiais na data da análise. Nenhuma generalização para outros postos ou
uma instalação real será feita sem evidência correspondente.
